\chapter{PIB Metric Bug and Plane Confusion}

\section{Why the Legacy Result Looked Exciting}

The \texttt{0330-focal-pib-sweep-4loss-cn2-5e14} run was the first attempt to optimize
focal-plane concentration explicitly. The figures looked impressive: the
\texttt{focal\_pib\_only} strategy produced a sharp focal peak and high PIB values across
multiple bucket radii. If interpreted naively, the run appeared to show that a
static D2NN had become an excellent detector preconditioner.

\section{Bug 1: PIB Is a Ratio}

The first problem is definitional. Power-in-bucket is a normalized fraction,
\begin{equation}
\PIB(r_b)=
\frac{\int_{|\rho|<r_b} I_f(\rho)\, d\rho}{\int I_f(\rho)\, d\rho}.
\end{equation}
By itself, a larger ratio does not guarantee a larger absolute detector signal.
The physically relevant quantity is the bucket power
\begin{equation}
P_{\mathrm{bucket}}(r_b)=\int_{|\rho|<r_b} I_f(\rho)\, d\rho
= \PIB(r_b)\,\TP\,E_{\mathrm{in}}.
\end{equation}
Here $E_{\mathrm{in}}$ is the energy incident on the diffractive block before
the detector lens, and $\TP$ is the output-plane throughput ratio
$\int |U_{\mathrm{out}}|^2 dA / \int |U_{\mathrm{in}}|^2 dA$. Later chapters use
both normalized received power and raw bucket energy, so this identity is the
bridge between those two tabulations.
If throughput collapses faster than the ratio improves, the detector receives
less power even though the normalized PIB appears better.

\section{Bug 2: Output Plane versus Focal Plane}

The second bug is geometric. Let
\begin{equation}
U_{\mathrm{out}} = H(U_{\mathrm{turb}}),
\end{equation}
and let the detector observe the focal-plane irradiance
\begin{equation}
I_f(\rho)=\left| \mathcal{L}\{U_{\mathrm{out}}\}(\rho)\right|^2.
\end{equation}
Then output-plane metrics such as $\CO$ and wavefront RMS live on
$U_{\mathrm{out}}$, while detector metrics such as $\PIB$ and Strehl live on
$I_f$. The legacy interpretation mixed these two planes, which made a detector
observable look like a field-recovery claim.

\runfigure{0330-focal-pib-sweep-4loss-cn2-5e14/focal_pib_only/06_fig1_focal_plane_vacuum_vs_turbulent_vs_d2nn.png}
  {Legacy focal-plane comparison}
  {The \texttt{0330} focal-plane result that initially suggested strong detector-side
  concentration.}

\runfigure{0330-focal-pib-sweep-4loss-cn2-5e14/focal_pib_only/07_fig2_pib_bar_chart_5um_10um_25um_50um.png}
  {Legacy PIB bars}
  {PIB increases across multiple bucket radii in the legacy interpretation.}

\runfigure{0330-focal-pib-sweep-4loss-cn2-5e14/focal_pib_only/08_fig3_d2nn_output_plane_irradiance_phase_residual.png}
  {Legacy output-plane field}
  {The output-plane field can look highly distorted even when the focal-plane
  bucket ratio appears favorable.}

\runfigure{0330-focal-pib-sweep-4loss-cn2-5e14/focal_pib_only/09_fig4_wavefront_rms_distribution_50samples.png}
  {Legacy wavefront distribution}
  {Wavefront diagnostics from the legacy run. These do not live on the same
  plane as the focal detector metric.}

\runfigure{0330-focal-pib-sweep-4loss-cn2-5e14/focal_pib_only/11_cheatsheet_d2nn_mode_conversion.png}
  {Legacy mechanism sketch}
  {Mechanism sketch retained as a qualitative reminder that static optics can
  redistribute energy without restoring the target field.}

\runfigure{0330-focal-pib-sweep-4loss-cn2-5e14/focal_pib_only/13_fig5_metrics_comparison_0330_vs_0401.png}
  {Legacy versus corrected metrics}
  {Comparison between the legacy interpretation and the corrected 0401
  re-evaluation.}

\section{How to Reclassify the 0330 Run}

The safe interpretation is therefore narrow:
\begin{itemize}
  \item the run is useful as a \emph{bug-discovery artifact};
  \item it is not a reliable benchmark for final detector performance;
  \item it motivated a corrected clean sweep with hardened data and explicit
  detector-side received-power accounting.
\end{itemize}

\begin{breakthrough}
The legacy PIB bug did not invalidate detector-plane optimization. It clarified
what the detector objective had to be: absolute received power on the focal
plane, not a normalized sharpness ratio interpreted in isolation.
\end{breakthrough}

That clarification forced a second cleanup step before retraining: the data and
physics pipeline also had to be rebuilt on a cleaner footing. That is the role
of the next chapter.
